\begin{lema}[Prescribed indices and long arrival]\label{lem:prescripcion-larga}
Let \(S_M\) be regular entry of the history at \(M=\lceil\delta n\rceil\) and \(H^0_j\) regular survival
up to and including \(j\) after that entry. If \(s\) rows and \(t\) columns are prescribed, with \(s,t\le
j/2\), the probability that they remain alive and \(H_j^0\) occurs, given the history at \(M\) on
\(S_M\), is at most
\[
 (j/M)^{s+t}\exp(C_{\rm pres}(st+t)/j).
\]
Under the condition on \(x_0\) in Lemma~\ref{lem:llegada-integrada}, \(\mathbb E_\nu[1+X;H^0_j\setminus
H_j]=O(e^{-2j})\), with \(X=Qe^{-L}\).
\end{lema}
\begin{proof}
If a mark is missing at \(M\), the probability is zero. At each regular parent of size \(k>j\), the
uniform row hits one of the \(s\) marks with probability \(s/k\). Given the row, the transition
probability is \(p_{il}(1+O(1/k))\) by the proof of Theorem~8.1 of~\cite{pc}, with class constants
independent of \(n/k\). The column sums of \(p\) equal one. Thus the probability of hitting one of the
\(t\) columns is \(\ge t/k-Ct/k^2\), and that of hitting both sets is \(\le Cst/k^2\) by diffusion.
Avoiding both sets of marks and remaining regular has probability at most
\((1-s/k)(1-t/k)\exp(C(st+t)/k^2)\), since the product is at least \(1/4\). Iterating gives
\(\prod_{k=j+1}^M(1-s/k)=(j)_s/(M)_s\le(j/M)^s\). The columns give the other factor, and
\(\sum_{k>j}k^{-2}\le1/j\).

On a diagonal there are at most \(M\) candidates at entry. If \(L\) sites remain, then for
\(r=\lfloor\alpha j/4\rfloor\) and \(R=\lceil\alpha j\rceil\) the indicator of \(L\ge R\) is bounded
above by \((L)_r/(R)_r\). Each ordered list requires \(r\) distinct rows and \(r\) distinct columns.
Since \((R)_r\ge(\alpha j/2)^r\), summing over fewer than \(4n\) diagonals gives
\[
 \nu(H^0_j\setminus H_j)\le
 4n\left(\frac{2j}{\alpha M}\right)^r
      \exp(C_{\rm pres}(r^2+r)/j).
\]
Dividing its logarithm by \(j\) gives a limit superior of at most \((\alpha/4)\log x_0+C_{\rm
pres}\alpha^2/16<-(\beta+3)\), since \(j=\lfloor bn\rfloor\) and \(b\) is fixed. From \(n_0(b)\) onward
the mass is \(\le e^{-(\beta+2)j}\). On regular arrivals, the static bounds of \cite[6.5--6.8]{tc} give
\(X\le e^{C_0+\beta j}\). For a board admitting no completion, \(X=0\). Multiplication by \(1+X\) gives
the stated cost. The margin \(3\) in the exponent covers rounding and the union bound factor before
absorbing \(\beta j\).
\end{proof}
